模型假设只保留会改变冲突判定、候选生成或跨问接口的规则，具体如下。

\begin{enumerate}[label=(\arabic*), leftmargin=*, itemsep=0.35em, topsep=0.25em, parsep=0pt]
  \item 区间与离散单位。时间区间和频段区间均采用半开区间，只有内部存在正长度交集才判定为重叠；原始端点及调整量按题面给出的离散单位取整数。该口径使首尾相接的事件不产生虚假冲突，并使连续区间检查能够与资源格复核逐格对应。\label{ass:interval}
  \item 重复事件与动作规则。一份计划的全部重复事件共享频段区间，频移或时移同时作用于该计划的全部事件；问题二中每份计划至多选择保留、频移、时移或撤销中的一种动作，并保持单次时长、间隔时长和使用次数不变。\label{ass:action}
  \item 跨问接口与 Q4 范围。问题二和问题三的基准资源域取 \([0,643)\times[0,100)\)，问题三冻结已经验证的问题二具体方案并采用频段宽 3、单次时长 2、间隔 8、使用 12 次的 C 类模板；问题四从问题一的 150 条原始计划独立起算，只允许 C 类计划取整数间隔 \(g_i'\in\{0,1,\ldots,18\}\)，不叠加频移或时移，也不读取 Q2 撤销和 Q3 新增计划。\label{ass:cross-interface}
\end{enumerate}

资源域、新增模板和问题四的间隔范围是为保证跨问一致性与可验证性而采用的明确建模口径；相关结果均在正文中标注其适用范围。
